% \chapter{Proposed Solution} \label{chap:proposed-solution}

% \section{Introduction}

% To address the limitations identified in existing \textit{IEC 61499} development environments, this dissertation proposes the design of a web-based \ac{IDE} for modelling \textit{IEC 61499}-compliant applications. The goal is to bridge the gap between modern collaborative software engineering practices and industrial automation tooling used in \ac{CPPS}.

% Existing tools are predominantly desktop-based, resource-intensive, and poorly suited for distributed development. In contrast, the proposed solution adopts a centralised web-based architecture that enables platform independence, simplified deployment, and real-time collaboration.

% \section{Approach}

% The approach is derived from functional and non-functional requirements extracted from the \textit{IEC 61499} standard and an analysis of existing tools such as \textit{4diac}. Functional requirements include support for function blocks, event and data connections, and deployment to compliant runtime environments such as \textit{DINASORE}.

% Non-functional requirements include:
% \begin{itemize}
%     \item Platform independence through browser-based access
%     \item Reduced setup complexity
%     \item Improved usability for distributed teams
%     \item Extensibility for future tooling enhancements
% \end{itemize}

% The proposed system is evaluated through the implementation of a prototype and qualitative comparison with existing desktop-based environments, focusing on usability, accessibility, and workflow efficiency.

% \section{System Architecture}

% The proposed system follows a centralised client--server architecture designed to support distributed \ac{CPPS} development.

% The system consists of three core components:
% \begin{itemize}
%     \item \textbf{Frontend \ac{IDE}}: A browser-based interface for modelling \textit{IEC 61499} applications.
%     \item \textbf{Server Middleware}: A centralised backend responsible for state management, collaboration, and communication.
%     \item \textbf{\acp{RTE}}: External execution environments such as \textit{DINASORE} where applications are deployed.
% \end{itemize}

% The server acts as an intermediary between developers and runtime environments, removing the need for direct device-level configuration on individual machines.

% \subsection{Communication Model}

% The system supports three main communication channels:
% \begin{itemize}
%     \item HTTP-based communication for serving the \ac{IDE}
%     \item Event-based communication for real-time collaboration
%     \item Socket-based communication for deployment to \acp{RTE}
% \end{itemize}

% This architecture enables centralised control while maintaining low-latency interaction between clients.

% \section{Real-Time Collaboration Concept}

% A key requirement of the proposed system is support for concurrent editing of structured \textit{IEC 61499} models.

% Traditional version control systems such as \textit{Git} are not well suited for graph-based representations due to frequent merge conflicts caused by structural changes rather than textual differences.

% To address this, the system adopts a conflict-free replicated data approach (\acp{CRDT}), enabling multiple users to edit shared application models simultaneously. Changes are propagated in real time and merged deterministically without conflicts.

% \section{Deployment Concept}

% Deployment in the proposed system is based on a centralised orchestration model. The server translates the graphical application model into a sequence of deployment commands that are executed on target \acp{RTE}.

% The deployment process follows these conceptual stages:
% \begin{enumerate}
%     \item Synchronisation of function block definitions across devices
%     \item Creation of function block instances on target \acp{RTE}
%     \item Establishment of event and data connections
%     \item Initialisation and execution start of the application
% \end{enumerate}

% This approach ensures consistent deployment across distributed runtime environments while abstracting low-level communication from the user.

% \section{Design Rationale}

% The choice of a web-based architecture is motivated by several factors:

% \begin{itemize}
%     \item \textbf{Accessibility}: Users can access the system without installing heavy desktop applications.
%     \item \textbf{Centralised management}: All configuration and deployment logic is handled by the server.
%     \item \textbf{Collaboration}: Real-time shared editing enables distributed teams to work concurrently.
%     \item \textbf{Scalability}: The architecture supports multiple users and distributed runtime environments.
% \end{itemize}

% Graph-based modelling is particularly suitable for \textit{IEC 61499}, as function block networks naturally map to node-based representations.

% \section{Methods}

% \subsection{Research Design}

% This research follows a design-oriented methodology, focusing on the development of a functional prototype. The aim is to demonstrate feasibility rather than statistical generalisation. The prototype serves as a research artifact used to explore architectural decisions and usability implications.

% \subsection{Evaluation}

% The system is evaluated through qualitative comparison with established \acp{IDE}, focusing on:
% \begin{itemize}
%     \item Usability
%     \item Accessibility
%     \item Feature completeness
%     \item Support for collaborative workflows
% \end{itemize}

% \chapter{Prototype Development} \label{chap:prototype-development}

% \section{Overview}

% This chapter describes the technical implementation of the proposed system. It details the server architecture, frontend implementation, collaboration system, and deployment mechanism used to realise the conceptual design presented in Chapter~\ref{chap:proposed-solution}.

% \section{Runtime Environment}
% The prototype targets \textit{DINASORE} (\textit{Dynamic INtelligent Architecture for Software and MOdular REconfiguration}) \cite{pereira_dinasore_2020} as the primary \ac{RTE}.

% \textit{DINASORE} was selected due to its open-source nature, simplified architecture, and suitability for prototyping. It exposes a command-based interface over port 61499, which is used by the \ac{IDE} for deployment communication.

% \section{Server Implementation}

% The server is implemented using \textit{Node.js}, providing a JavaScript-based runtime that supports both backend logic and compatibility with frontend libraries.

% The server is responsible for:
% \begin{itemize}
%     \item Serving the web-based \ac{IDE}
%     \item Managing application state and storage
%     \item Handling real-time collaboration updates
%     \item Orchestrating deployment to \acp{RTE}
% \end{itemize}

% Four main communication types are handled:
% \begin{itemize}
%     \item HTTP requests for static content
%     \item WebSocket communication for collaboration
%     \item HTTP POST requests for deployment triggers
%     \item Socket communication with \acp{RTE}
% \end{itemize}

\section{Frontend Implementation}

The frontend provides the graphical interface for modelling \textit{IEC 61499} applications.

\subsection{Devices Configuration View}

This view allows configuration of distributed runtime devices. Each device contains:
\begin{itemize}
    \item Name and address
    \item \textit{DINASORE} port
    \item SSH configuration
    \item Visual colour identifier
\end{itemize}

This centralised configuration removes the need for per-machine setup.

\subsection{Function Block Editor}

Function blocks are defined using:
\begin{itemize}
    \item XML-based interface definitions (\texttt{.fbt})
    \item Python-based behaviour implementation
\end{itemize}

This separation follows the \textit{IEC 61499} encapsulation model.

\subsection{Graph View}

The Graph View is the main modelling environment where function blocks are connected into executable systems.

A node-based editor is used to represent function blocks and their connections. Each connection enforces type compatibility based on \textit{IEC 61131-3} data types.

The prototype uses a graph editor library selected for its balance between simplicity and flexibility, supporting typed inputs, multi-port nodes, and real-time updates.

\subsection{Console View}

The Console View displays runtime logs and deployment output, providing visibility into system execution and debugging.

\section{Real-Time Collaboration Implementation}

Real-time collaboration is implemented using \ac{CRDT}-based data structures.

The system uses shared state replication to ensure consistency across clients without merge conflicts.

The graph is stored as two shared maps:
\begin{itemize}
    \item Node map (function blocks, positions, mappings)
    \item Link map (connections between ports)
\end{itemize}

Each update is propagated in real time and merged deterministically across all connected clients.

\section{Deployment Implementation}

\subsection{Function Block Synchronisation}

Before deployment, function block files (\texttt{.fbt} and \texttt{.py}) are synchronised to all target devices via SSH using an automated script.

\subsection{Command Construction}

Deployment commands are encoded into a binary format required by \textit{DINASORE}. Each message consists of:

\begin{itemize}
    \item A configuration name header
    \item A payload header
    \item The command data
\end{itemize}

\subsection{Deployment Sequence}

Deployment follows a strict ordered sequence:
\begin{enumerate}
    \item Query runtime environment
    \item Create execution resource (\texttt{EMB\_RES})
    \item Create function block instances
    \item Write initial values
    \item Establish connections
    \item Start execution
\end{enumerate}

Each request requires confirmation before the next is issued.

\subsection{Runtime Monitoring}

After deployment, the system issues watch requests to monitor function block outputs. Values are continuously retrieved and displayed in the Graph View for debugging and validation.

\section{Console and Debugging}

Deployment logs and runtime messages are displayed in the Console View, including:
\begin{itemize}
    \item Synchronisation status
    \item Command execution traces
    \item Runtime responses from \acp{RTE}
\end{itemize}

\section{Summary}

The prototype demonstrates the feasibility of a web-based collaborative \textit{IEC 61499} development environment, integrating modelling, deployment, and runtime monitoring into a single unified system.
